% Generated draft: requires mathematical and source review.
\documentclass[11pt]{article}
\usepackage[margin=25mm]{geometry}
\usepackage{fontspec}
\setmainfont{Latin Modern Roman}
\usepackage{amsmath,amssymb,mathtools,unicode-math}
\setmathfont{Latin Modern Math}
\usepackage{enumitem,needspace,xurl,hyperref}
\setlength{\parindent}{0pt}
\setlength{\parskip}{5pt}
\setlength{\emergencystretch}{3em}
\newcommand{\N}{\mathbb{N}}
\newcommand{\Z}{\mathbb{Z}}
\newcommand{\Q}{\mathbb{Q}}
\newcommand{\R}{\mathbb{R}}
\newcommand{\C}{\mathbb{C}}
\newcommand{\F}{\mathbb{F}}
\newcommand{\A}{\mathbb{A}}
\newcommand{\PP}{\mathbb{P}}
\newcommand{\E}{\mathbb{E}}
\newcommand{\eps}{\varepsilon}
\newcommand{\dd}{\,\mathrm{d}}
\begin{document}
\section*{UnsolvedMath ID: 30004560}
\textbf{Source problem code:} catalogue entry
\subsection{Definitions and mathematical statement}
Let \(G\) be a weighted bipartite graph with sphere topology and a marked face \(f\). Thus the vertices have a bipartition \(V(G)=B\sqcup W\), every edge joins the two parts, and the graph is considered with an embedding in the sphere and a distinguished face. Write \(w\) for its edge weights and \(G^*\) for the augmented dual intended in the source. The source requires \((G,w)\) to be ``suitably non-degenerate,'' but does not supply the precise admissibility conditions, the domain of the weights, or the construction of the augmented dual.

For a source-admissible triple \((G,w,f)\), let \(\mathcal P(G,w,f)\) denote the collection of p-embeddings associated with perfect t-embeddings of this augmented dual, in the intended source terminology. This notation is a placeholder: the supplied record does not define either embedding notion or the correspondence between them. Let \(\mathbb R^{2+1}\) denote Minkowski space and let \(\mathbb H\) denote the unit hyperboloid referred to in the source. Write
\[
\operatorname{Iso}_{\mathrm{Lor}}(\mathbb R^{2+1};\mathbb H)
=
\{L: L\text{ is a Lorentz isometry of }\mathbb R^{2+1},\ L(\mathbb H)=\mathbb H\}.
\]
The precise hyperboloid and any sheet or orientation restrictions are not specified in the record.

The source asks whether every suitably non-degenerate weighted bipartite graph with sphere topology and a marked face has an augmented dual admitting a perfect t-embedding, described in the clean statement as equivalently admitting a p-embedding. It further asks whether the associated p-embedding is unique up to the above Lorentz isometries. In the placeholder notation, the intended existence and uniqueness assertions are
\[
\forall\,(G,w,f)\ \text{source-admissible},\qquad
\mathcal P(G,w,f)\ne\varnothing,
\]
and
\[
\forall\,(G,w,f)\ \text{source-admissible},\quad
\forall P_1,P_2\in\mathcal P(G,w,f),\quad
\exists L\in\operatorname{Iso}_{\mathrm{Lor}}(\mathbb R^{2+1};\mathbb H)
\quad P_2=L\cdot P_1,
\]
where \(L\cdot P_1\) denotes the intended geometric action on the p-embedding. These formulas express the requested target conditionally on the missing definitions; they do not supply those definitions or establish the asserted equivalence of embedding notions. A fully specified mathematical statement requires the source's admissibility hypotheses, embedding definitions, boundary conditions, and precise equivalence relation.
\subsection{Short English statement}
The problem asks whether every suitably non-degenerate weighted bipartite graph embedded in the sphere, with a marked face, admits the source's perfect t-embedding of its augmented dual, equivalently its associated p-embedding, and whether that p-embedding is unique up to Lorentz isometries preserving the specified unit hyperboloid. The supplied record omits definitions and hypotheses needed to make this a fully precise assertion.
\subsection{Sources}
\begin{itemize}
\item Local UnsolvedMath registry, supplied record with permanent ID \(30004560\), title ``Existence and Uniqueness of Perfect T-Embeddings,'' exact source problem code \texttt{catalogue entry}, in the Oberwolfach Reports Open Problems collection. No individual proposer is identified; the supplied proposed year is 2020.
\item Supplied report citation: \emph{Mini-Workshop: Dimers, Ising and Spanning Trees beyond the Critical Isoradial Case} (2021), \emph{Oberwolfach Reports}, pp.\ 1759--1789, DOI \texttt{10.4171/owr/2020/35}. \url{https://doi.org/10.4171/owr/2020/35}. The record attributes the problem to workshop open-problem discussions but supplies no problem-specific page or speaker attribution.
\item Chelkak, Laslier and Russkikh, ``s-embeddings'' (supplied reference label): \url{https://arxiv.org/abs/2109.06272}.
\item Berggren, Nicoletti and Russkikh, ``p-embeddings'' (supplied reference label): \url{https://arxiv.org/abs/2408.05441}.
\item Galashin, ``amplituhedra and origami'' (supplied reference label): \url{https://arxiv.org/abs/2410.09574}.
\item Literature assessment dated exactly \texttt{2026-08-22}: supplied dated curation, not independently verified here. The supplied status is ``partially solved.'' The assessment states that p-embedding existence is established in important special families and related t-embedding existence results are available, while general existence and the required perfect or canonical p-embedding theory remain unresolved. The record expressly requests independent verification of time-sensitive status and cited claims.
\item The registry labels the statement \texttt{corrected\_verified} and reports that it was checked against the official Oberwolfach report, preserving the report's hypotheses, notation, and open target. This is a supplied verification claim, not an independent verification in this draft.
\end{itemize}
\end{document}
